\section{Discussion}\label{sec:discussion}

\textbf{Why the approach works, and where it does not.} The decisive step is the formulation:
self-referencing the restriction to each instance's own baseline turns wells of very different pressure
into one learnable problem, and the survival objective uses the many censored windows instead of
mislabelling them. Every learned forecaster built on this formulation outperforms the physics rule and the
anomaly detector. Within the deep models, the recurrent architectures rank and alarm best: the informative
signal is a slow trend in two or three channels, which a recurrent state integrates directly, whereas a
patch Transformer splits it into 16-minute tokens and must re-assemble it through attention from only 24
positive events. The physics branch, the mask-aware channel-independent encoder and the ensemble give
PI-PatchSurv properties that matter in operation, namely calibrated intervals, an accurate time-to-event
and robustness to failed sensors, but not a better ranking.

\textbf{Practical implications.} With a threshold fixed on other wells, the best forecasters warn of
roughly two thirds of the events a median of 2--4~h before the restriction reaches $+5\,\%$, at 1.2--1.6
false alarms per well-month, and they raise no alarm on normal operation or on the quick-restriction
look-alike. Hours of warning are enough to schedule a choke cleaning within a shift but not to plan an
inhibitor campaign; longer warnings require longer real histories than 3W provides. Self-referencing the
score to each well's own recent level is a simple, model-agnostic step that roughly doubles deployed recall
for several models and should be part of any deployment.

\textbf{Strengths.} The study evaluates on wells excluded from training with all thresholds and
hyper-parameters chosen on separate calibration wells; it compares twelve models with fixed and tuned
hyper-parameters and several seeds; it reports calibration, intervals, event-level and confounder-level
false alarms, robustness and a dataset-version generalisation test; and every number is regenerated from
saved predictions by published code.

\textbf{Failure cases.} (i) On well 24, PI-PatchSurv's score drifts within long instances and produces
false alarms under the self-referenced threshold. (ii) Hydrate plugs downstream of the choke (classes 8 and
9) raise the upstream pressure slowly and are confused with scale; separating them needs the missing
downstream sensors or temperature dynamics. (iii) Wells 6 and 21 contribute few events and are often
missed. (iv) The autoencoder detector does not forecast at all.

\textbf{Threats to validity.} \emph{Internal:} only 24 real events are available and 17 come from one well,
so fold-level variance is large and the deep models cannot be separated statistically; ablations other than
three components use a single seed. \emph{Construct:} $\RI$ is a proxy based on one pressure and does not
control for changes of the choke opening, which is missing in most instances; the $+5\,\%$ level is a
modelling choice. \emph{External:} 3W covers offshore wells of one operator with short recordings;
results may differ for onshore gas wells or other scale chemistries. \emph{Conclusion validity:}
early stopping and tuning rely on a single small calibration well per fold, which makes both noisy.

\textbf{Were the research gap and objectives met?} The first gap, the absence of timed forecasting of
scaling, is addressed: calibrated multi-horizon forecasts on unseen wells are shown to be feasible, but the
lead time is hours, not the days targeted in the proposal, because of the data. The second gap, the lack of
fusion of chemistry and sensors, is addressed only partly: physical indicators are fused with learned
features and help, while the chemistry prior from public brines helps little. The third gap, the absence of
unseen-well evaluation, is fully addressed. The fourth gap, rare events and missing uncertainty, is
addressed by censor-aware training, synthetic data and conformal intervals with near-nominal coverage. The
objective of a proposed model that outperforms all baselines was not achieved; the evidence instead
supports the formulation and the evaluation protocol, with simple recurrent forecasters as the strongest
rankers.
